% Crescent Moon: one-pass object detection from a small vision-language model.
% Build from paper/: latexmk -pdf main.tex
% Figures are read from ../artifacts/figures/ (scripts/figures.py output).
\documentclass[10pt,twocolumn]{article}
\usepackage[letterpaper,top=1in,bottom=1in,left=0.82in,right=0.82in,columnsep=0.3in]{geometry}
\usepackage[T1]{fontenc}
\usepackage{mathptmx}
\usepackage{microtype}
\usepackage{graphicx}
\usepackage{booktabs}
\usepackage{amsmath}
\usepackage[numbers,sort&compress]{natbib}
\usepackage[hidelinks]{hyperref}
\usepackage{xcolor}
\usepackage{caption}
\usepackage{titlesec}
\usepackage{tikz}
\usetikzlibrary{positioning,calc,arrows.meta}
\captionsetup{font=small,labelfont=bf,skip=6pt}
\titleformat{\section}{\large\bfseries}{\thesection}{0.6em}{}
\titleformat{\subsection}{\normalsize\bfseries}{\thesubsection}{0.6em}{}
\titlespacing*{\section}{0pt}{1.4ex plus 0.4ex}{0.7ex}
\titlespacing*{\subsection}{0pt}{1.1ex plus 0.3ex}{0.4ex}
\setlength{\bibsep}{2pt plus 0.3pt}
\graphicspath{{../artifacts/figures/}}
\renewcommand{\bibfont}{\footnotesize}
\newcommand{\ok}{\textcolor{green!45!black}{+}}
\newcommand{\no}{\textcolor{red!70!black}{$-$}}
\newcommand{\eq}{$\circ$}
\newcommand{\pend}{$\ast$}

\title{\LARGE\bfseries Crescent Moon: One Pass Object Detection\\ from a Small Vision-Language Model}
\author{Jakob Frick\\ \normalsize\texttt{jakob.maria.frick@gmail.com}}
\date{}

\begin{document}
\twocolumn[{%
\maketitle
\vspace{-1.2em}
\begin{center}
\includegraphics[width=0.86\textwidth]{speed.jpg}
\captionof{figure}{Accuracy against time per image on COCO (left) and Pascal VOC (right),
on one L40S. Slots match Moondream~3.1 in less than half its time.}
\label{fig:speed}
\end{center}
\vspace{0.6em}
}]

\begin{abstract}
Small vision-language models (VLMs) have learned through their pretraining a deep
understanding of the world, but when used for object detection they have to generate
predictions one token at a time one prompt per category. Inspired by diffusion-based
drafting models for text-based LLMs, we implement object detection as a one pass task. We
extend a frozen Moondream~2 model (1.9B) with slots that a single forward pass fills to
predict object positions using the model's existing coordinate, size and vocabulary heads
read. We post-train low-rank adapters at the slot positions and an objectness layer with a
Hungarian set loss on 46k COCO images, and 17.7k teacher-labelled images.

We find that our model matches the Moondream~3.1 (9.3B) F1 score while taking less than
half the time (0.091\,s vs 0.208\,s) on COCO val2017 images. On Pascal VOC 2007, which no
model was tuned on, it scores 0.858 in 0.063\,s against 0.857 in 0.095\,s, suggesting that
the model is using the general image understanding of the base model. While being slower
than traditional, much smaller object detection models such as YOLO11x, we compare them on
COCO-O, to demonstrate the robustness of the VLM-based approach. We find that our model
can detect objects continuously across frames in streamed video, while maintaining
quality.

Our approach finds fewer objects smaller than Moondream~3.1 (recall 0.20 against 0.36),
which we were able to close by implementing reading of a $2\times2$ grid of tiles at the
cost of increased latency. This demonstrates this approach allows for a tradeoff between
speed and accuracy. We show ablation about different parameters and give a log of
different approaches that were implemented.
\end{abstract}

\section{Introduction}
\begin{figure*}[t]
\centering
\definecolor{slotA}{HTML}{2A78D6}\definecolor{slotB}{HTML}{EB6834}
\definecolor{slotC}{HTML}{1BAF7A}\definecolor{slotD}{HTML}{6250D6}
\definecolor{frozenfill}{HTML}{EEF3FA}\definecolor{trained}{HTML}{E8833A}
\resizebox{\textwidth}{!}{%
\begin{tikzpicture}[
  font=\footnotesize, x=1cm, y=1cm,
  tok/.style={minimum size=0.26cm, inner sep=0, draw=black!40, fill=black!10},
  empty/.style={minimum size=0.26cm, inner sep=0, draw=black!45, dashed, fill=white},
  full/.style={minimum size=0.26cm, inner sep=0, draw=#1!75!black, fill=#1!85},
  none/.style={minimum size=0.26cm, inner sep=0, draw=black!25, fill=black!4},
  lab/.style={font=\scriptsize, text=black!65, align=center},
  head/.style={font=\footnotesize\bfseries, align=center},
  arr/.style={-{Latex[length=1.8mm]}, black!60, thick},
]
% input image
\node[inner sep=0] (in) at (0,0) {\includegraphics[height=2.7cm]{architecture-input.jpg}};
\node[head] at (0,1.85) {Image};
% input sequence: image tokens, then empty slots
\node[head] at (3.75,1.85) {Input};
\foreach \i in {0,...,5} \node[tok] at (2.6+\i*0.3,1.05) {};
\node[lab, anchor=west] at (4.3,1.05) {image tokens};
\foreach \r in {0,...,3} \foreach \i in {0,...,5} \node[empty] at (2.6+\i*0.3,0.45-\r*0.4) {};
\node[lab] at (3.35,-1.25) {48 empty slots};
\draw[arr] (1.95,0) -- (2.3,0);
% the model
\draw[draw=blue!35!black!45, fill=frozenfill, rounded corners=3pt] (6.0,-1.35) rectangle (8.9,1.35);
\node[align=center] at (7.45,0.5) {Moondream 2\\[1pt]\scriptsize frozen};
\draw[draw=trained, fill=trained!22, rounded corners=1.5pt] (6.2,-1.1) rectangle (8.7,-0.45);
\node[font=\scriptsize, text=trained!65!black, align=center] at (7.45,-0.775) {adapters at the slots\\trained};
\node[head] at (7.45,1.85) {One forward pass};
\draw[arr] (4.5,0) -- (6.0,0);
\draw[arr] (8.9,0) -- (9.35,0);
% output: the same slots, filled
\node[head] at (10.9,1.85) {Output};
\foreach \c/\t/\r in {slotA/person/0,slotB/horse/1,slotC/person/2}
  { \foreach \i in {0,...,5} \node[full=\c] at (9.65+\i*0.3,0.85-\r*0.4) {};
    \node[anchor=west, font=\scriptsize] at (11.5,0.85-\r*0.4) {\t}; }
\foreach \i in {0,...,5} \node[none] at (9.65+\i*0.3,-0.35) {};
\node[anchor=west, font=\scriptsize, text=black!55] at (11.5,-0.35) {nothing};
\node[font=\scriptsize] at (10.4,-0.75) {$\vdots$};
\node[lab] at (10.4,-1.3) {x \,y \,size \,+\, label};
% detections
\node[inner sep=0] (out) at (14.6,0) {\includegraphics[height=2.7cm]{architecture-output.jpg}};
\node[head] at (14.6,1.85) {Detections};
\draw[arr] (12.5,0) -- (out.west);
\end{tikzpicture}}
\caption{We append empty slots to the image tokens of a frozen Moondream~2. One forward
pass fills every slot with a box and a label, or leaves it as ``nothing''.}
\label{fig:architecture}
\end{figure*}

Object detection has been one of the traditional applications of Machine Learning for
over a decade~\citep{girshick2014rcnn,redmon2016yolo}. Vision Language Models also
localise objects through text by emitting coordinate tokens~\citep{chen2022pix2seq}, but
decoding is sequential and conditioned by prompts. Moondream~2's
\texttt{detect}~\citep{moondream} takes three decoder steps per box and one call per
phrase. Latency grows with categories and objects in the image. This makes VLMs unfit for
fast object detection especially when there is no prompt given.

Set prediction (DETR~\citep{carion2020detr}) uses a fixed number of queries who each
return a box or ``no object'' in one pass to perform object detection. It is the basis of
today's most accurate detectors, DETR and its descendants such as
DINO~\citep{zhang2023dino} and of real-time variants such as
RT-DETR~\citep{zhao2024rtdetr}.

Inspired by diffusion-based drafting models in modern language
models~\citep{austin2021d3pm,sahoo2024mdlm,nie2025llada}, we ask whether a diffusion
approach would allow a VLM to perform fast object detection.

We append 48 slots to a frozen Moondream~2 and train about 50M adapter parameters
(Section~2). We find that one pass is enough and that additional diffusion passes do not
improve detection quality. The resulting head is level with Moondream~3.1~\citep{moondream31}, a 9.3B model, on COCO and VOC
in 2.3 and 1.5 times less time (Section~3, Fig.~\ref{fig:speed}). Its time does not
depend on the number of categories, which makes detection on streamed video practical
(Table~\ref{tab:video}). It finds fewer very small objects unless it also reads tiles,
and it knows only COCO's 80 categories (Section~4). The aim is reliable rather than
exhaustive detection.





\section{Method}

\subsection{Slots}
Moondream~2 turns an image into 730 image tokens, which its language model then reads to
produce text. We append 48 \emph{slots} after these tokens (Fig.~\ref{fig:architecture}).
A slot is a group of six positions that starts empty and can hold one object: three
positions for the box (its centre x, centre y, and size) and three for the label, which
is long enough for every COCO category name.

One forward pass of the model fills all slots at once. We do not add new output layers
to read them. Moondream~2 already has a head that turns a position into a coordinate or
a size, and a head that turns a position into a word, and we use both as they are. The
only new output is an \emph{objectness layer}: a single linear layer that looks at each
slot and outputs one number, the probability that the slot holds a real object. Slots
with a low probability are treated as ``nothing''.

This is the first step of how a masked diffusion model
writes~\citep{chang2022maskgit,austin2021d3pm,sahoo2024mdlm,nie2025llada}: every position
starts as a blank and all of them are predicted together. Such a model would then keep
its most confident predictions and predict the rest again, for several steps. We stop
after the first.

\subsection{Training}
The weights of Moondream~2 stay frozen. We train three things: the starting embeddings
of the slots, the objectness layer, and low-rank adapters (LoRA~\citep{hu2022lora})
that are only active at the slot positions. Because the adapters do not touch the image
tokens, the model still captions and answers questions exactly as before. At rank 64 this
is about 50M trained parameters, and about 100M at rank 128.

The model has to learn which slot should describe which object, and no slot is assigned
in advance. We follow DETR~\citep{carion2020detr}: for every training image, the
Hungarian algorithm~\citep{kuhn1955hungarian} pairs each labelled object with the slot
whose prediction is closest to it. A paired slot is trained to produce that object's box
and label, with a cross-entropy loss on the coordinates and the label and an L1 and
GIoU loss~\citep{rezatofighi2019giou} on the box. All other slots are trained to say
``nothing''. A run takes two to four hours on one H100 (batch size 16, 20{,}000 steps).

\subsection{From slots to detections}
Each slot gets a score: its objectness times the confidence of its label and its box. We
keep the slots above a threshold. When two kept boxes with the same label overlap by an
IoU of 0.5 or more, we merge them into one~\citep{solovyev2021wbf}.

Small objects are hard to find, because the vision encoder shrinks the image to a grid
of $27\times27$ cells. As an option, the model can also read a $2\times2$ grid of
overlapping crops of the image (\emph{tiles}) in the same batch and add the small boxes
it finds there~\citep{akyon2022sahi}. This finds more small objects and takes 2.5 times
as long. We choose the threshold and these settings on 500 COCO validation images that
are not part of the test set.

\subsection{Training data}
We train on 46k images from COCO train2017~\citep{lin2014coco}, with its 80 categories.
To see more kinds of images, we add images from Open
Images~\citep{kuznetsova2020openimages} and Objects365~\citep{shao2019objects365} that a
larger VLM labelled for us, a form of distillation~\citep{hinton2015distilling}. To pick
this \emph{teacher}, we compared twelve VLMs through OpenRouter~\citep{openrouter} on 200
COCO images. Gemini~4 Carbon~\citep{gemini2023} agreed best with COCO's labels (F1 0.756)
at \$0.005 per image. It labelled 21k images for \$120, and 17.7k of them passed our
label filter.

\section{Experiments}
\begin{table*}[t]
\centering\small\setlength{\tabcolsep}{7pt}
\caption{F1 at IoU 0.5 on the four test sets, and seconds per COCO image on one L40S.}
\label{tab:main}
\begin{tabular}{@{}lccccr@{}}
\toprule
Model & COCO & VOC & COCO-O & Generated & s/img \\
\midrule
Moondream 2 \texttt{detect} & .586 & .851 & .773 & .641 & .220 \\
Moondream 3.1 & .700 & .857 & \textbf{.781} & \textbf{.795} & .208 \\
YOLO11n & .612 & .770 & .410 & .582 & \textbf{.008} \\
YOLO11x & .715 & .838 & .520 & .734 & .018 \\
\midrule
Slots, rank 64 & .705 & .840 & .765 & .779 & .091 \\
Slots, rank 64, with tiles & .714 & .841 & .773 & \textbf{.796} & .226 \\
Slots, rank 128 & .707 & \textbf{.858} & .772 & .786 & .091 \\
Slots, rank 128, with tiles & \textbf{.718} & .844 & .775 & .791 & .229 \\
\bottomrule
\end{tabular}
\end{table*}


\subsection{Setup}
We compare the slot model with Moondream~2, the model it is built on, with
Moondream~3.1~\citep{moondream31}, a newer and five times larger model from the same
family, and with YOLO11~\citep{yolo11}, a dedicated detector trained on COCO, in its
smallest (n, 2.6M parameters) and largest (x, 57M) size.

We use four test sets. \emph{COCO} is 500 images from val2017. \emph{VOC} is 1{,}000
images from the Pascal VOC 2007 test set, which none of our choices were tuned on.
\emph{COCO-O}~\citep{mao2023cocoo} is 6{,}782 human-labelled images of COCO's categories
that do not look like COCO's photos: sketches, cartoons, paintings, handmade objects,
tattoos and photos in bad weather. \emph{Generated} is a set we made ourselves: 3{,}950
images generated with a text-to-image model in 88 photographic and artistic styles, and
labelled by the teacher. No model saw these images in training, but their labels come
from the same teacher that labelled part of our training data.

We report F1 at an IoU of 0.5, where each predicted box can match one labelled box.
Moondream has to be asked for one category at a time, so we ask it for every category
that is present in the image. The slot model and YOLO are not prompted. All times are
measured on one L40S GPU.

\subsection{Accuracy}
The slot model is as accurate as Moondream~3.1 (Table~\ref{tab:main}). At rank 128 it
scores 0.707 on COCO against 0.700, and 0.858 on VOC against 0.857. Differences of this
size are within noise: a paired bootstrap puts the noise at about 1.5 points, so we call
this a tie. Compared with Moondream~2, the model it is built on, it gains 12 points on
COCO.

The two models reach this score in different ways. The slot model is more careful: on
COCO its precision is 0.811 against 0.682, and its recall is 0.626 against 0.719. Most of
what it misses is small. For objects smaller than one cell of the feature grid its
recall is 0.20, where Moondream~3.1 reaches 0.36. Reading tiles closes this gap (0.37)
and raises COCO F1 to 0.718, at 2.5 times the time. On VOC, where objects are larger,
tiles do not help.

Rank 64 and rank 128 score the same on COCO, and rank 128 is 1.8 points better on VOC.
Both take the same time, so we use rank 128 in the rest of this section.
Table~\ref{tab:coco} in the appendix gives precision and recall for every model.

\begin{figure*}[t]
\centering
\includegraphics[width=0.86\textwidth]{selection-precision-recall.jpg}
\caption{Precision against recall on COCO (left) and VOC (right). The slot model gives a
confidence per box, so it has a curve; the dot is the threshold we use. Moondream and the
teacher give no confidence and are single points.}
\label{fig:pr}
\end{figure*}

\subsection{Speed}
\begin{figure}[t]
\centering
\includegraphics[width=0.9\columnwidth]{asked.jpg}
\caption{Time per image against the number of categories asked.}
\label{fig:asked}
\end{figure}
The slot model takes 0.091\,s per COCO image and Moondream~3.1 takes 0.208\,s
(Fig.~\ref{fig:speed}). Filling the slots is a small part of that: 0.016\,s. The rest is
the vision encoder and the image tokens, which Moondream needs as well.

The gap grows with the number of categories. Moondream runs once per category, and the
slot model runs once per image. When we ask for all 80 COCO categories, the slot model
still takes 0.08\,s, Moondream~2 takes 1.52\,s and Moondream~3.1 takes 2.67\,s
(Fig.~\ref{fig:asked}). Asking for everything also hurts Moondream's accuracy: it returns
a box for 71\% of the categories that are not in the image.

YOLO is much faster than any of these models, at 0.008 to 0.018\,s per image.

\subsection{Robustness}
On COCO, the largest YOLO matches the slot model in a fifth of the time (0.715 against
0.707). If the images look like COCO and only COCO's categories are needed, a dedicated
detector is the better choice.

This changes when the images look different. On COCO-O the slot model scores 0.772 and
stays within a point of Moondream~3.1 (0.781), while YOLO11x falls to 0.520 and YOLO11n
to 0.410. YOLO still places its boxes well, but it finds only half as many objects. On
our generated images the gap is smaller: the slot model scores 0.786 and YOLO11x 0.734.
We read this as the slot model keeping the general image understanding of the VLM it is
built on, which a detector trained only on COCO does not have.

Figure~\ref{fig:pr} shows where this comes from in training. Adding the teacher-labelled
images moves the slot model's curve on COCO towards Moondream~3.1, and the teacher itself
(F1 0.728 on COCO) is still ahead of every model we tested.

\subsection{Video}
Because one pass covers every category, the slot model can run on video as it arrives.
At 0.10\,s per image, detecting every third frame keeps up with 30 frames per second.

We tested this on two street scenes from MOT17~\citep{milan2016mot16}, looking for
people in every third frame. The slot model scores F1 0.447 at 0.104\,s per frame, and
Moondream~3.1 scores 0.429 at 0.206\,s (Table~\ref{tab:video} in the appendix). With
tiles it finds the most people (F1 0.496).

Detecting each frame on its own makes boxes flicker: 10 to 18\% of the people found in
one frame are missed in the next. A simple tracker that links boxes between frames, and
keeps a box for one frame after it disappears, lowers this to 9 to 12\% without changing
the F1 or the time.

\subsection{What mattered}
\emph{One step is enough.} We expected that filling the slots would need several
refinement steps, as in a diffusion model. It does not. Decoding in two or three steps
gave the same F1 as one. We think this is because an object is only six positions, and
each one is read by a head the model already has, which leaves little to refine. We
have not tested this explanation.

\emph{Three changes account for most of the accuracy.} The first is reading boxes
through Moondream's own coordinate head instead of predicting four numbers directly
(F1 0.429 to 0.563). The second is data: four times more images and longer training
(to 0.622), keeping small objects in the labels (0.636), and adding the teacher-labelled
images (0.657 to 0.673). The third is how we pick the final boxes, which added 2.7 points
without any training. Table~\ref{tab:ablate} in the appendix lists every change we
tried, including those that did not help.

\emph{Smaller weights.} A phone or browser cannot hold the full model, so we tested it
with every weight of Moondream~2 rounded to 4 bits, which is a 1.1\,GB download instead
of 3.9\,GB. With one image crop, the rank-64 model then scores 0.654 on COCO instead of
0.705, and 0.721 on COCO-O, which is still 20 points above YOLO11x.

\section{Limitations}
The slot model only knows COCO's 80 categories, while Moondream can be asked for anything
in free text. We have not tested it on larger vocabularies such as
LVIS~\citep{gupta2019lvis} or on referring expressions~\citep{yu2016refcoco}. It is
trained on COCO images, and Moondream has probably seen COCO as well, so VOC is the
fairer comparison of the two. It misses more small objects than Moondream unless it reads
tiles. All times come from one type of GPU, and our COCO and VOC test sets are samples of
500 and 1{,}000 images.

\section{Conclusion}
We asked whether a small vision-language model could detect objects the way a diffusion
language model drafts text, and found that a single drafting step is enough. A frozen
Moondream~2 (1.9B) extended with 48 slots and a post-trained adapter of 50--100M
parameters returns every object of every category in one forward pass. It matches the F1
of Moondream~3.1 (9.3B) on COCO and on Pascal VOC, while being significantly faster and
maintaining static inference time regardless of the number of categories. It maintains
the robustness of the backbone it uses: on the shifted images of COCO-O it stays within a
point of Moondream~3.1, while a dedicated detector with the same COCO accuracy loses a
quarter of its F1. The most promising use we see for the model is continuously on
streamed video, because of its static inference time that is below 100\,ms.

\section{Future work}
The most direct extension is to carry slots across video frames. The denoising pass
already trains the head to complete a set of boxes of which some positions are revealed;
giving it the previous frame's boxes and labels as those positions would let each slot
keep its object, so that the model tracks as it detects instead of linking boxes
afterwards, and such a model should be evaluated on full sequences with identity-aware
measures such as HOTA rather than per-frame F1. A second direction is cost per frame: the
vision encoder dominates the time of a pass, so reusing its features where the scene has
not changed, tiling only where small candidates appear, or replacing it with a cheaper
encoder would each bring the model closer to dedicated detectors, and would remove the
2.5-fold cost that small objects carry today. Finally, the head is limited to COCO's 80
categories. Open-vocabulary labels are longer than three tokens, and it is there that
decoding in more than one step, which did not help here, may begin to matter.

\bibliographystyle{plainnat}\bibliography{references}

\clearpage
\onecolumn
\appendix
\section{Every change we tried}

Table~\ref{tab:ablate} gives each change by itself, with COCO F1 before and after. Rows
marked \ok{} helped, \eq{} stayed within noise (about 1.5 points by paired bootstrap), and
\no{} hurt. Three rows that left COCO unchanged moved something else: human Objects365
and Open Images labels in place of the teacher's lifted VOC F1 by 1.3 points, tile crops
of small objects in training lifted it by 1.7 and raised precision but not small-object
recall, and a size-weighted loss raised recall under one cell by 0.03. On the same 22k
images, adding the teacher's extra boxes to COCO's raised tiled small-object recall from
0.37 to 0.41 and lost precision. Mean-teacher
self-training~\citep{tarvainen2017meanteacher,liu2021unbiasedteacher,xie2020noisystudent}
lost 0.9 points. Separately from the head, custom fused CUDA kernels for Moondream~2's
sequential decoder cut its cost per box 1.85 times; their output is close to eager
PyTorch's but not bit-identical.

\begin{table}[h]
\centering\small\setlength{\tabcolsep}{3pt}
\caption{COCO F1 at IoU 0.5 before and after each change.}
\label{tab:ablate}
\begin{tabular}{@{}clr@{}}
\toprule
& Change & F1$_{.5}$ \\
\midrule
\ok & Coordinate bins instead of regressed boxes (11k img.) & .429$\rightarrow$.563 \\
\ok & 4$\times$ data, 3$\times$ steps & .563$\rightarrow$.622 \\
\ok & Keep small objects (small recall .10$\rightarrow$.26) & .622$\rightarrow$.636 \\
\ok & LoRA rank 16$\rightarrow$64 & .636$\rightarrow$.657 \\
\ok & + 18k teacher-labelled images & .657$\rightarrow$.673 \\
\ok & Combined score, duplicate removal & .674$\rightarrow$.694 \\
\ok & Fusion at IoU 0.5, box confidence$^{1/4}$ & .700$\rightarrow$.707 \\
\ok & $2\times2$ tiles (recall $<$1 cell .20$\rightarrow$.37) & .707$\rightarrow$.718 \\
\eq & LoRA rank 64$\rightarrow$128 & .694$\rightarrow$.700 \\
\eq & 2--3 decoding passes instead of 1 & unchanged \\
\eq & No denoising pass in training & .705$\rightarrow$.705 \\
\eq & Category prompts before the slots & .705$\rightarrow$.702 \\
\eq & Size-weighted loss (recall $<$1 cell +.03) & .705$\rightarrow$.700 \\
\eq & Human O365/OI labels for teacher's (VOC +.013) & .705$\rightarrow$.698 \\
\eq & Small-object tile crops in training (VOC +.017) & .705$\rightarrow$.696 \\
\no & Two slot groups by object size (23k img.) & .590$\rightarrow$.549 \\
\no & Objectness target = matched IoU (small recall $-$.07) & .636$\rightarrow$.631 \\
\no & ``Nothing'' only near labelled objects & .636$\rightarrow$.568 \\
\no & Teacher labels only, no COCO (17.7k vs.\ 46k img.) & .694$\rightarrow$.681 \\
\no & Teacher's labels for COCO's, same 22k img. & .692$\rightarrow$.678 \\
\no & + teacher boxes COCO lacks, same 22k img. & .692$\rightarrow$.678 \\
\no & EMA self-training, ignore confident slots & .694$\rightarrow$.685 \\
\bottomrule
\end{tabular}
\end{table}

\section{Video}

\begin{table}[h]
\centering\small\setlength{\tabcolsep}{3.5pt}
\caption{People in two MOT17 street scenes, every third frame.}
\label{tab:video}
\begin{tabular}{@{}lcccr@{}}
\toprule
Model & P & R & F1$_{.5}$ & s/frame \\
\midrule
Moondream 2 \texttt{detect} & .638 & .308 & .415 & .260 \\
Moondream 3.1 & .816 & .291 & .429 & .206 \\
Slots r128, fast & .652 & .341 & .447 & .104 \\
Slots r128, tiled & .559 & \textbf{.446} & \textbf{.496} & .418 \\
Slots r64, COCO only, fast & \textbf{.851} & .241 & .376 & \textbf{.087} \\
\bottomrule
\end{tabular}
\end{table}

\section{COCO in detail}\label{app:coco}

\begin{table}[h]
\centering\small\setlength{\tabcolsep}{5pt}
\caption{Precision, recall and F1 on COCO for the models of Table~\ref{tab:main}.}
\label{tab:coco}
\begin{tabular}{@{}lccccc@{}}
\toprule
Model & Precision & Recall & F1$_{.5}$ & F1$_{.75}$ & Recall $<$1 cell \\
\midrule
Moondream 2 \texttt{detect} & .501 & .706 & .586 & .450 & -- \\
Moondream 3.1 & .682 & \textbf{.719} & .700 & .560 & .36 \\
Slots, rank 64 & .779 & .644 & .705 & .555 & .23 \\
Slots, rank 64, with tiles & .763 & .670 & .714 & \textbf{.589} & .33 \\
Slots, rank 128 & \textbf{.811} & .626 & .707 & .566 & .20 \\
Slots, rank 128, with tiles & .743 & .695 & \textbf{.718} & .586 & \textbf{.37} \\
\bottomrule
\end{tabular}
\end{table}

\section{Example detections and areas under the precision-recall curve}

Examples are the first images whose annotation count falls in a set range and that
every model's evaluation covers; boxes are drawn as each model returned them, with the
slot head at the threshold tuned on COCO. Captions give boxes and seconds per image.

\begin{figure}[h]
\centering
\includegraphics[width=\textwidth]{examples.jpg}
\caption{COCO val2017.}
\end{figure}

\begin{figure}[h]
\centering
\includegraphics[width=\textwidth]{examples-voc.jpg}
\caption{Pascal VOC 2007 test.}
\end{figure}

\begin{figure}[h]
\centering
\includegraphics[width=\textwidth]{examples-mot.jpg}
\caption{MOT17-09, frames 1, 4 and 7 (every third frame), people only.}
\label{fig:mot}
\end{figure}

\begin{figure}[h]
\centering
\includegraphics[width=0.92\textwidth]{selection-areas.jpg}
\caption{Area under the precision-recall curve for every slot model we trained, without tiles (solid) and with tiles (hatched).}
\label{fig:areas}
\end{figure}
\end{document}
